\paragraph{Discrepant analysis in diagnostic medicine.} Using a new test to decide which specimens get the reference test, and keeping the old result for the rest, was common practice in the evaluation of nucleic-acid amplification tests in the 1990s.

\emph{The bias.} \citet{hadgu1996discrepancy,hadgu1997bias,hadgu1999discrepant} and \citet{miller1998bias} showed that the procedure biases the new test's sensitivity and specificity upwards. \citet{green1998evaluation} and \citet{lipman1998quantifying} quantified the size of this bias, and \citet{hadgu2000discrepant} reviewed the debate.

\emph{The fixes.} The composite reference standard of \citet{alonzo1999using} and the analysis of \citet{hawkins2001some} restore validity by also sampling concordant results, which is our two-phase design in the clinical setting. Related lines of work cover verification bias \citep{begg1983assessment} and latent-class models for imperfect reference standards \citep{hui1980estimating,walter1988estimation,reitsma2009review}. \citet{limmathurotsakul2012fools} illustrate how an imperfect gold standard can mislead.

\emph{Our contribution relative to this work.} The bias itself is therefore not new. What we add is the analysis that matters for machine learning:
\begin{itemize}
\item an exact identity for any number of models and classes, rather than one binary test;
\item its consequences for rankings and for the flagger, which is itself an evaluated model;
\item lock-in under panels of correlated flaggers;
\item sharp identified sets under the discrepant design;
\item and an optimality theory for the corrective design.
\end{itemize}

\paragraph{Pooling bias in information retrieval.} Test-collection construction in information retrieval faces the same structure. Only documents retrieved by participating systems are judged, and unjudged documents are assumed non-relevant. This favours systems that contributed to the pool and penalises new ones \citep{zobel1998how,buckley2007bias}, which is the retrieval counterpart of our lock-in result. Estimators based on probability sampling of relevance judgments with inverse-probability weights \citep{aslam2006statistical,yilmaz2006estimating} are the retrieval counterpart of our two-phase design. We are not aware of work that carries this experience over to the relabelling of classification and question-answering benchmarks, where the pooled ``systems'' are the evaluated models themselves.

\paragraph{Label errors in test sets.} Confident learning \citep{northcutt2021confident} and model-based annotation-error detection \citep{klie2023annotation,chong2022detecting} made it routine to find label errors with models, and \citet{northcutt2021pervasive} showed that such errors can change benchmark rankings. Several authors note in passing that flag-only review leaves errors unaccounted for. \citet{northcutt2021pervasive} describe their noise-prevalence estimates as ``optimistic in favor of higher capacity models''; \citet{vendrow2025do} call their error counts lower bounds; and \citet{chong2022detecting} warn that cleaning biases benchmarks towards models similar to the cleaner. Section~\ref{sec:practice} documents how widespread the design is. Correlated errors among language models \citep{kim2025correlated} explain why hidden errors are frequent when both the labels and the flagger come from such models. The same correlation is a known problem for aggregating LLM judges \citep{balasubramanian2026dependence} and goes back to the conditional-independence assumption of Dawid--Skene \citep{dawid1979maximum}. \citet{lam2003evaluating} and \citet{chavoshi2026impact} bound the effect of a noisy reference on estimated accuracy. They do not consider an adjudication design in which part of the benchmark is verified, which is what makes our bounds sharp.

\paragraph{Label-efficient and prediction-powered evaluation.} Our estimator is the classical difference estimator \citep{cassel1976some,sarndal1992model}, a form of double sampling \citep{tenenbein1970double,neyman1938contribution}, with the original label as the auxiliary variable. In the language of prediction-powered inference it is PPI \citep{angelopoulos2023prediction,angelopoulos2026ppi} with the original label as the prediction and the flag as a stratifier \citep{fisch2024stratified}, or active inference with flag-dependent inclusion probabilities \citep{zrnic2024active}. \citet{mozer2026ppi} makes the equivalence between PPI and the difference estimator explicit. We claim no novelty for the estimator.

Recent work on evaluating models with few labels uses an autorater's verdict on a model's outputs as the proxy \citep{boyeau2025autoeval,chatzi2024prediction,eyre2025regression,fogliato2024framework,chen2026efficient}. Active testing selects informative test points \citep{sawade2010active,kossen2021active}. In our setting the proxy is the benchmark label that is already there, and the sampling strata are defined by the very disagreement that motivates the audit. The relevant baseline is not a cheaper estimator but a biased design in wide use.
